% ============================================================
%  Poul Martin Møller,
%  "Forberedelser til en Afhandling om Affectation"
%
%  A fragment (Brudstykke IV): preparations toward a treatise on
%  affectation. Per the editor's title-note, Møller had long
%  collected remarks on this "side of life" (cf. the Strøtanker),
%  intending a separate work; in 1837 he began a connected
%  scientific treatment, printed here first, followed by the
%  editor's thematic grouping of the scattered observations.
%
%  Published posthumously in:
%  Efterladte Skrifter af Poul M. Møller, bd. 3 (Tredie Bind),
%  ed. F.C. Olsen. Kjøbenhavn: C.A. Reitzel.
%  (Transcribed from the 3rd edition, "Tredie Udgave".)
%  Printed pp. 163--188. PDF (bibliotek scan) pp. 174--199;
%  offset PDF = printed + 11.
%
%  TRANSCRIPTION (Danish)
%  Conventions: original orthography kept (aa-spellings,
%  capitalised nouns, æ/ø, "philosophisk", "Charakteer",
%  "Reflexion", "speculativt"). Fraktur long-s -> s. Danish
%  quotes „…" kept. Letterspaced emphasis -> \emph{}. Printed
%  page numbers -> \opage{N} (marginal, at the point the printed
%  page begins). Editor's section-groupings marked by a centred
%  rule (\ombreak).
% ============================================================

\documentclass[12pt]{article}
\usepackage[utf8]{inputenc}
\usepackage[T1]{fontenc}
\usepackage{graphicx}
\DeclareUnicodeCharacter{0254}{\reflectbox{c}}
\usepackage{libertinus}
\usepackage{libertinust1math}
\usepackage{textalpha}
\usepackage[danish]{babel}
\usepackage[protrusion=true,expansion=true]{microtype}
\usepackage[margin=1.4in]{geometry}
\usepackage{setspace}
\usepackage{fancyhdr}
\usepackage{marginnote}
\usepackage[colorlinks=true,linkcolor=black,urlcolor=black,
  pdftitle={Forberedelser til en Afhandling om Affectation},
  pdfauthor={Poul Martin Møller}]{hyperref}
\setlength{\headheight}{15pt}
\pagestyle{fancy}
\fancyhf{}
\fancyhead[LE,RO]{\thepage}
\fancyhead[RE]{\textit{Om Affectation}}
\fancyhead[LO]{\textit{Poul Martin Møller}}
\setlength{\parindent}{1.5em}
\setlength{\parskip}{0pt}
\onehalfspacing

\newcommand{\opage}[1]{\marginnote{\footnotesize\textit{[#1]}}}
\newcommand{\ombreak}{\par\medskip\begin{center}\rule{0.32\textwidth}{0.4pt}\end{center}\medskip}

\title{\textbf{Forberedelser til en Afhandling\\[0.2em] om Affectation}\footnote{Længe havde Poul Møller havt Opmærksomheden henvendt paa denne Side af Livet og nedskrevet sine Bemærkninger derover, hvortil da knyttede sig den Hensigt at bearbeide dette Emne i et eget Skrift. Man vil ved Læsningen af de foregaaende „Strøtanker`` have bemærket denne staaende Retning i hans Iagttagelse, saa at der næsten paa hver Side findes Træk af Falskhed og Selvbedrag i Menneskenes Liv. Blandt disse vare ogsaa de Optegnelser, som her følge særskilt, og hvilke fornemmelig ere saadanne, som Forfatteren selv havde betegnet som nedskrevne med Hensyn til hiint Begreb. Denne Adskillelse er vilkaarlig, da Meget af det, der er blevet staaende i Samlingen, lige saavel hører herhid; men dersom Alt, hvad der kunde henføres til Begrebet Affectation, var blevet borttaget af Hovedsamlingen, saa vilde denne have tabt sin eiendommelige, paa en Maade biographiske, Charakteer, og nogen Vilkaarlighed vilde dog endnu blive tilbage ved Bestemmelsen af den Grændse, hvortil det efter egen Fortolkning af de nedskrevne Tankers Tendents var tilladt at gaae. Desuden har Forfatteren kun løselig angivet Maaden, hvorpaa han vilde have forbundet denne talrige Mængde af Bemærkninger og specielle Træk, og langtfra samlet dem for at udfylde et forud betænkt Schema, men især stræbt at fremdrage Affectationen i dens mere latente Yttringer; saa at disse Iagttagelser, saaledes som de haves fra Forfatterens Haand, i ethvert Tilfælde dog kun vilde virke temmelig sporadisk paa Læseren. De, som nu følge, ere af Udgiveren sammenstillede partiviis efter deres Indhold. I Aaret 1837 begyndte Poul Møller paa en sammenhængende videnskabelig Behandling af denne Gjenstand, og hvad han saaledes har nedskrevet, er her meddeelt først.}}
\author{Poul Martin Møller}
\date{Efterladte Skrifter, bd.\ 3 (Brudstykke IV)}

\begin{document}
\maketitle

\begin{center}\rule{0.32\textwidth}{0.4pt}\end{center}
\medskip

\opage{163}I de her meddeelte Reflexioner har Forfatteren vel søgt at tydeliggjøre Affectationens Begreb, men det er --- som en ikke \opage{164}ganske flygtig Læser selv let bliver vaer --- ikke det, som har udgjort hans Hovedinteresse. Det har meest været ham om at gjøre, at gjennemgaae Affectationens almindeligste Phænomener, saavel i Menneskelivet overhovedet som i Særdeleshed i Tidsalderen. Den, som ikke finder Smag i dette Slags moralske Naturbeskrivelser, advares derfor betimeligen om at springe disse Sider forbi.

De, der ikke sætte nogen Priis paa anden Tænkning end den, hvori alle Begreber ved en immanent Udvikling frembringes af Intet, ville naturligviis fordre, at det her afhandlede Begreb skal fremkomme som Resultat af et forudskikket speculativt System. Men hvis disse ikke, i Følge den alt givne Advarsel, have hørt op at følge med os, anmodes de om at frafalde deres Fordring, da det her ansees for upassende, at give Grundtegningen til et Tempel, for deri at finde et Hul for en fattig Kirkerotte. Udgjøre virkelig alle deres theoretiske og praktiske Begreber eet System, da ville de sikkert deri have bestemt det her omhandlede Begreb med de samme Grændser, som i det Følgende skulle antydes. --- Den anden Classe Læsere, hvis enkelte Begreber hver for sig føre en mere selvstændig Tilværelse og mere have ordnet sig i Formen af et Archipelagus, ville sandsynligviis i een eller anden Gruppe finde Plads for disse Reflexioner. I modsat Fald kunne de vel tilveiebringe den Smule Rystelse i deres Tankers Ocean, hvorved en liden Holm kan komme frem og bestaae for en halv Times Tid.

Adskillige dybsindige Mennesker i vor Tid betjene dem i Skrifter og Samtale af Ordet Sandhed, som de synes at tage i en saadan Betydning, at de derved betegne næsten Alt, hvad der er agtværdigt i Menneskelivet. De fordømme da paa den anden Side som Usandhed eller vel endogsaa indre Løgn saamange Slags fornuftstridende Sindsstemninger, at man naturligviis bringes til at opkaste det Spørgsmaal, om Sandhed hos \opage{165}dem blot er en af Sædelighedens væsentlige Sider, eller om alle Sædelighedens Bestemmelser ere forenede der, eller om det er en særegen Dyd ved Siden af andre Dyder? Dette synes de Mange, der i aphoristiske Yttringer betjene sig af hiint Begreb, ei ret at have gjort sig rede for; og der gives neppe noget systematisk Skrift, hvori hiin Synsmaade saaledes er gjennemført, at det ubestemte Begreb, som foresvæver hine, deri kan antages at have sin adæqvate Fremstilling. Af de korte Registre, man undertiden finder over de forskjellige Moralprinciper, der have været gjorte gjeldende i den praktiske Philosophies Historie, kunde man fristes til at troe, at Engellænderen Wollaston netop havde opfattet den moralske Fuldkommenhed som Sandhed; men ved at læse hans eget Værk finder man ikke deri den Synsmaade, som her tales om. Heller ikke findes nogen Udførelse deraf i Ammons Moralsystem, uagtet man ved at eftersee hans Kapitel om Moralprincipet vil bringes til at troe det.

Et Slags Sandhed har et Menneskes Liv, naar det uden Forstillelse følger sin naturlige Begjering: en høiere Sandhed har det, naar det har opnaaet Dyden (i den antike Betydning af Ordet), saa at det endnu af de naturlige Drifter henter sine Handlingers Indhold, men dog har vundet et saadant Herredømme over dem, at det i deres Tilfredsstillelse iagttager et vist Maadehold. Et endnu høiere Trin af personlig Sandhed har dens Liv, der med reen fornuftig Autonomie bestemmer alle sine Forsætter. (At dette ikke kan skee ved en blot subjectiv Tænkning, men saaledes, at Subjectet erkjender den fornuftige Orden, som uden dets Medvirkning er i sin Udvikling i Tilværelsen, for at være et Værk af den samme Fornuft, der er dets Villies egentlige Sandhed, hører det ikke herhid at bevise.) For saavidt endelig Menneskets rene Selvbestemmelse er den ved Religiøsitet helliggjorte Villie, handler det i fuldkommen Harmonie med den hele Fornuftverden; det er, hvad det skal være, og dets Liv kan \opage{166}ei opnaae nogen høiere Sandhed. Men denne Sandhed er ikke andet end Sædelighed, og al Afvigelse fra den er Usædelighed. Affectationen er vist nok en Art af denne Usandhed; men vi have nærmere at bestemme, hvad Plads den indtager i denne større Kreds.

Det Menneske, der har Affectation i sit Liv, bestemmer sig forsaavidt ikke med fuldkommen moralsk Frihed, dets Handlinger have ei deres Udspring af det sande Selv, som er dets frie sædelige Villie. Dets Villie bestemmes af et eller andet blot naturligt Øiemed, hvorved det bringes til at forestille en fremmed Person eller paatage sig en falsk Rolle, som ei er det anviist i Livet.

Det hører da for det Første til Affectationens Væsen, at den er en Falskhed. Men ikke enhver Falskhed er Affectation. Den, der med klar Bevidsthed om, at han lyver eller forstiller sig, gjør sig skyldig i Løgn og Forstillelse, lægger ikke deri nogen Affectation for Dagen. Vi sige hermed ingenlunde, at Affectation er nogen ringere Last end den Forstillelse, der er ledsaget med klar Bevidsthed. Det vil i det Følgende vise sig, at en Grad af Affectation forudsætter mindre Immoralitet end den planmæssige Forstillelse, medens der i en anden Grad deraf lægges større Immoralitet for Dagen. Affectation er da ikke ublandet Falskhed, men har altid en Tilsætning af Selvbedrag; thi det ligger i Affectationens Begreb, at Personen stræber at være, hvad han ikke kan være; men derefter kan han ikke stræbe, uden i det Mindste for en Stund at indbilde sig selv, at han kan være det. Men det aldeles uforskyldte Selvbedrag, som er uden al Falskhed, kan ligesaa lidet kaldes Affectation; saa at de to her angivne Bestanddele: Falskheden og Selvbedraget, stedse findes forbundne deri. (For at opnaae et udenfor Handlingen liggende Øiemed.)

(Det burde være en overflødig Bemærkning, at aldeles ufri\opage{167}villige Bevægelser umulig kunne henregnes til Affectation, og at følgelig de abnorme Livsyttringer, der formedelst det aandelige Livs Mangel paa Raadighed over det legemlige, ikke sjelden undslippe Mennesket aldeles mod dets Villie, maae udelukkes derfra. Det er dog ikke sjeldent, at man hører saadanne unaturlige og uadæqvate Yttringer af det indre Liv belægges med Navnet Affectation. Vi tænke herved for det Første paa alle de udvortes Bevægelser, der paa Grund af organiske Feil ikke staae i det Forhold til Personens Forestillinger, Tanker og Følelser, hvor de vilde staae, hvis Legemet var i sin normale Tilstand. Exempler herpaa ere keitet Bevægelse af enkelte Lemmer, ufrivillige Muskeltrækninger, Feil i Udtalen \&c., hvilke tilsammen af Pøbelen henregnes til Affectation. Under samme Rubrik henhøre endeel forseilede Livsyttringer, der ere af ligesaa uskyldig Art, skjøndt selv fornuftige Folk ikke sjelden afsige Fordømmelsesdom over dem. Dertil kan man regne enhver uforholdsmæssig Yttring af Følelser, der bliver uforholdsmæssig, fordi Personen enten bestandig eller under visse Omstændigheder mangler Evne til at finde det træffende Udtryk for sin Stemning. Saadanne Feiltagelser kunne ligesaa lidet henregnes til moralske Feil som Udlændingens Ord, naar han af Mangel paa Sprogfærdighed kommer til at betjene sig af uhøviske eller fornærmende Udladelser.\footnote{I et ældre Udkast staaer her følgende: Betragter man nemlig ethvert udvortes Meddelelsesmiddel som Sprog --- og i denne Sammenhæng ere jo Miner og Bevægelser ligesaa vel et Sprog som det lydeligen udtalte Ord --- da gives der en stor Deel ufrivillige Sprogfeil, der uden Grund regnes til Affectation.} En Mand af ringe Stand, der ikke er vant til at befinde sig i Selskab med fornemme Folk, gjør ofte for dybe Buk; ikke fordi han vil lægge saa stor Ærefrygt for Dagen, men fordi han ikke har den sædvanlige Færdighed i at beregne Inclinationsvinklen \opage{168}med Nøiagtighed. --- Den, der under en Conversation bemærker, at han mod sin Villie kommer til at yttre sig for haardt og stærkt mod en Anden, vil undertiden gjøre sin Forseelse god igjen, men yttrer, fordi hans Blod nu engang er kommet ud af Ligevægten, en Velvillie, der ei er saa gandske meent, skjøndt den forulykkede Correction dog kan have sit Udspring af en Stræben efter at lægge sit sande Hjertelag for Dagen.\footnote{I et ældre Udkast: Det behøver vel ei at bemærkes, at vi ved Skildringen af disse Begreber gaae ud fra den ueftergivelige Forudsætning, at Vedkommende ei har andet Øiemed end at yttre, hvad der virkelig boer i ham. Er der andre Bihensigter med i Spillet, som f.\ Ex.\ hos den, der, fordi han vil skjule sine Tankers langsomme Circulation og sin Mangel paa Aandsnærværelse, griber en feil Yttring af sin Personlighed, netop fordi han griber til paa Slump, da finder naturligviis ingen fuld Oprigtighed Sted, skjøndt Andre just ikke kunne bedømme, i hvad Grad den savnes.} Alt i Menneskets Livsyttringer, som saaledes slet ikke ligger i Subjectets forudgaaende Forestillinger om dem, kan ikke kaldes Affectation. Mennesket ligger i slige Tilfælde kun under for en Strid med Naturens Nødvendighed, idet udvortes Organer vægre sig ved at adlyde hans frie Villies Befalinger. (Grund i foregaaende Affectation\footnote{I et ældre Udkast: Men hermed skal det ingenlunde være sagt, at det her omtalte Misforhold mellem Personligheden og dens Yttring ikke kan have sin Grund i foregaaende Affectation, (et oprindeligen paataget Væsen, som dens, der behager sig i Andres aldeles individuelle Legemsbevægelser, deres Gang eller Stemme, fordi enten disse Yttringer for sig betragtede forekomme ham at klæde denne godt, eller fordi dennes udmærkede Personlighed saaledes fortryller og forblinder ham, at den for ham kaster en Glands paa saadanne tilfældige Smaating. Saaledes fortæller Plutarch, at Alexander den Stores Hofmænd fik den Vane, at lade deres Hoved hælde noget til den ene Side, fordi de bestandig saae deres Fyrste i denne Stilling. Saadanne Livsyttringer, der først ere udsprungne af Affectation, kunne tilsidst være blevne til en anden Natur, saa at de ikke længer fremkaldes af Forsæt eller ledsages af Bevidsthed.)}).) Den paatagne Rolle kan være Efterligning, der \opage{169}strider mod Ens Natur eller en efter eget Behag digtet Charakteer, som synes at klæde En.)

Betragte vi Affectationens Grader efter dens større eller mindre Sammenhæng med Subjectets Charakteer, da har samme den mindst mulige Deel deri, naar den strider mod Villiens Retning, men i det enkelte Moment undslipper Mennesket, fordi dets Dyd ei er bleven til Færdighed. En saadan Yttring af Affectation staaer i samme Forhold til Affectationen, der er bleven til Vane, som en Beruselse har til Drikfældighed, og der gives neppe nogen Dødelig, som ikke i det Mindste i Ungdommen undertiden gjør sig skyldig i smaae, halv ufrivillige Nærligheder af dette Slags. Gives der nogen, som endogsaa i Begyndelsen af sin Udviklingsperiode saaledes bevarer sin Selvstændighed, at han endog undgaaer den Skikkelse af Affectationen, der ligger nær ved Elskværdighedens Grændse, at han f.\ Ex.\ aldrig, af Lyst til at være fuldkommen enig med en Ven eller Veninde, tvinger sig til at sympathisere mere med dem, end han kan, da maa han vist nok af Naturen have en meget stærkt udpræget Individualitet; men om han derfor har det heldigste Naturel, er et andet Spørgsmaal. Hans Gave til at bevare sin Individualitet reen og fri for uvedkommende Tilsætninger, kan ogsaa have sin Grund i et egoistisk Hang til at indeslutte sig blot i sin egen Tankekreds, og en Mangel paa Evne til at aabne sit Sind for fremmed Indflydelse. Den, der ei er i Stand til i Omgang at give sig saaledes hen til Andre, at han for en Stund bliver til Eet med dem, at han gandske gaaer ud af sig selv og taber sig i en fremmed Bevidsthedskreds, han kan vist nok ved sin Tilbageholdenhed redde sig for at overvældes af nogen aandelig Magt; men den \opage{170}Individualitet, der kun paa denne Maade bjerges, bliver altid eensidig og fattig.

Den aandelige Fuldkommenhed kan ligesom den physiske Væxt kuns fremmes derved, at Individet jævnlig smelter sammen med, hvad der er fremmed for det, og tilsyneladende opoffrer sig selv for at vende beriget hjem igjen til sig selv. Men da kan det vist nok skee, at en velbegavet Yngling fra en saadan frugtbringende Selvforglemmelse, hvori han har ladet en Anden raade for sin Tankegang, vender tilbage til sig selv med Skinnet af Forestillinger og Følelser, som forskydes af hans sande Jeg, men som han paatvinger sig selv for at nyde Sympathiens henrivende Lyst. De paa den Maade fremkomne Yttringer af Affectation kunne tidt netop blive Grund til en Bestyrkelse af det personlige Livs Sandhed, idet den, der af Vanvare har overskridt sin Personligheds Grændse, af Utilfredshed med sig selv bringes til klarere Bevidsthed om hvad han føler, vil og veed, og til en kraftigere Selvbevægelse i lignende Tilfælde. Den, der har været paa hiin Side Grændsen, veed nu først Besked om, hvor Grændsen er. --- En saadan forbigaaende Yttring af Affectation, der har sin Kilde i en Følelsen forblindende Kraft, fra hvilken Mennesket ved en paafølgende Reflexion løsriver sig, som fra Noget, der ikke hører til dets Væsen, og som det for Fremtiden vogter sig for, er en mindre Forseelse end den egentlige Løgn. (\emph{Den momentane Affectation.})

En høiere Grad af Affectation finder Sted hos den, der har antaget Vanen til et bestemt Slags falske Yttringer, idet han indbilder sig at have visse Meninger, Interesser eller Tilbøieligheder, fordi han af en eller anden udvortes Grund ønsker at have dem. Som naar En af Forfængelighed tillyver sig Kjærlighed for en eller anden Kunst, som han ikke har nogen Sands for, eller naar Slægtninge og Venner af en Zelot for dem selv og Andre gjelde for at være ligesindede med ham, uagtet hans Iver er deres Hjerte fremmed og under forandrede Omstændig\opage{171}heder af sig selv falder bort. Muligheden af Selvbedraget beroer her paa den logiske Conseqvents, hvormed den paatagne Rolle forfølges; men det Centrum, der bærer en saadan Forestillingskreds, falder udenfor Subjectet selv. Denne anden Grad af Affectation er i moralsk Henseende slet ikke mindre tilregnelig end den af klar Bevidsthed ledsagede Løgn og Forstillelse; thi at den ikke bringes til fuld Bevidsthed, kommer blot af, at den Paagjeldende ei vil bringe den til Bevidsthed. Dette bemærkes, fordi et høist urigtigt Begreb herom er kommet i Gang, i det Mængden troer, at der i Selvbedraget ligger tilstrækkelig Retfærdiggjørelse, eller i det Mindste Undskyldning endogsaa for aabenbar Slethed, --- en Mening, hvis Urigtighed let indlyser, naar man forfølger den i nogle af dens unegtelige Conseqventser. Deraf vilde nemlig følge, at man skulde kunne drive sin Falskhed saa vidt, at den igjen blev til Ærlighed, og at den, der lyver saa længe, til han selv troer sine digtede Fortællinger, paa dette Trin skulde blive bedre, fordi hans Løgn forvandles til Selvbedrag. Efter denne Synsmaade kunde en Dommer ogsaa i al Uskyldighed lade sig bestikke, da han blot behøvede ved Sophismer at bringe sig selv til at troe, at hans Velgjører havde Ret. Den Overbeviisning, der har sit Udspring af Tilbøielighed, kan ikke være synderlig grundig, hvilket allerede viser sig deraf, at Virkningen i Almindelighed ikke varer længer end dens Aarsag. Det er ikke usandsynligt, at den, der finder et literært Arbeide slet, fordi han troer, at hans Fjende har skrevet det, vil finde det meget godt, dersom han erfarer, at hans Ven er Forfatter dertil. Men at han i første Tilfælde fandt det slet, var vist nok ikke nogen uskyldig Vildfarelse, om der end var aldrig saa meget Selvbedrag deri; thi man bør ligesaa lidet i sine Antagelser som i sine Handlinger følge sine Tilbøieligheder imod Fornuft og Sandhed. Under denne Affectationens anden Grad optager Mennesket et falsk Element i sig og forvansker sin Personlighed, saa \opage{172}at dens Yttringer ikke hænge sammen med dets egentlige Selv.

For saavidt det saaledes har et dobbelt Indre, eet virkeligt, som er tilbagetrængt, og eet tilsyneladende, som det vil skal gjelde for det selv og Andre, fører det med Hensyn til det sidste kuns et Skinliv. (\emph{Den faste Affectation.})

Den sidst omtalte Affectation, selv hvor den er dreven allervidest, tillader dog, at Personen har varigt Skin af indvortes Sammenhæng saavel for Subjectet selv som for Andre. Men saadant Skin af Conseqvents kan kun være momentant i den tredie og værste Grad af Affectation. Det finder Sted, hvor Mennesket ikke har et eller andet skrømtet Træk i sin Charakteer, hvor det ikke har Bane til et bestemt Slags Affectation, men har en Færdighed i Affectation i Almindelighed, som snart antager denne, snart hiin bestemte Skikkelse. Denne Slethed nærmer sig efter sin større eller mindre Udvikling meer eller mindre til fuldkommen Usandhed i det personlige Liv. Hvis et Menneske kunde naae dens Culminationspunct, da vilde der ikke være nogen blivende Kjerne i dets Tænken og Villen, men det dannede sig i ethvert Øieblik af dets Liv en temporær Personlighed, for at ophæve den i det næste. Det vilde vel for det Meste, som visse Dyr, skifte Farve efter dets Omgivelser, og forsaavidt være det passive Product af dets Forhold; men da dette kuns er een af Affectationens Skikkelser, vilde dets Bane dog ikke kunne lade sig beregne efter den simple Regel, at det skulde ligne sine Omgivelser, da Affectation ogsaa kan vise sig i en Stræben efter i sin Vandel at fremstille det Særegne, det Usædvanlige. Til denne fuldendte Løgn i det indre Liv kan Ingen bringe det; men hvis Nogen kunde det, da vilde det være et moralsk Selvmord, hvorved han aldeles havde tilintetgjort sig som særegen Figur i den moralske Verden. (\emph{Den vexlende Affectation.})

Efter at vi i det Foregaaende have udviklet Affectationens tre Hovedarter med Hensyn til den større eller mindre Sammen\opage{173}hæng, den har med Charakteren, skulle vi nu betragte de forskjellige Maader, hvorpaa den yttrer sig, og de forskjellige Grunde, hvoraf den har sit Udspring.

\ombreak

Affectation kan vise sig deri, at man med halv Forsæt efterligner det Særegne ved Andres ufrivillige Livsyttringer, fordi det synes at klæde godt. Det er et Slags Affectation, som undertiden er en reen Frugt af Kjærlighed og Beundring uden andet Øiemed. Det Paatagne, Fremmede i de Bevægelser, som skulle være ufrivillige, og som Ingen kan stræbe at forandre, uden der er nogen helbredelig Naturfeil deri, kaldes af Alle meest Affectation, skjøndt det er det ubetydeligste Slags Affectation. --- Efterligne det Tilfældige i Andres Sprog. Bruge philosophisk Terminologie uden den philosophiske Indsigt.

Den kan lægge sig for Dagen i skrømtet Følelse. Menneskene sætte sig da saa levende ind i deres Sted, som virkelig have en saadan Følelse, at de bilde sig selv ind, de have den. Det kan man gjøre: det viser jo varme Skuespilleres Exempel, da de saaledes kunne identificere sig med en Rolle, at de glemme, hvem de ere. --- Det gjøre Leilighedsdigtere \emph{sensu pejore}, der undertiden sætte sig ind i sørgmodig Stemning, undertiden i glædelig, alt eftersom man har bestilt et Sørgedigt eller en Bryllupsvise hos dem. At Folk kunne bilde sig selv ind, at de have Følelse, kommer altsaa af, at de bruge Følelsens Sprog og Udtryk, og frembringe derved et Slags Reminiscents af den sande Følelse hos sig. --- Manden, der talte sig i Sorg over den af\opage{174}døde Slægtning.\footnote{Jvfr.\ S.\ 34.}

Man kan paatage sig en Selvfølelse, som man ikke har, ligesaa godt som en social. Ogsaa den høiere contemplative Følelse lader sig eftergjøre ved en Kunst, der skuffer baade Andre og de Affecterede selv. Især er Begeistring for Kunst en sædvanlig Affectation. (Tilløiet Enthusiasme --- Lystighed.)

Affectation kan vise sig i Erkjendelsen, naar En uden Overbeviisning af Tilbøielighed antager Sætninger for sande.

Affectationen viser sig ogsaa deri, at Mennesket under Selvbedrag yttrer sig, som om det havde en Tilbøielighed, et Naturkald, som det ikke har. --- (Kunstner.)

\ombreak

Affectationen har altid sin Grund i at Mennesket er bestukket af en eller anden Tilbøielighed, uden selv at vide det. Thi den Affecterede gjør efter Forsæt, hvad der ikke har nogen Betydning, uden hvor det er fremkaldt ved Naturmagt og Naturtilskyndelse eller høiere Nødvendighed. Denne Tilbøielighed er snart Egennytte, snart Svaghed.

Saaledes bringer Egennytten mangen Overtroens Præst til at antage Troeslærdomme, fordi hans Stilling i Livet beroer paa, at de ere sande. En katholsk Geistlig f.\ Ex.\ maa forkynde mangen en Lærdom, som i sig selv er saa stridende mod Fornuften, at man ikke kan antage, at saa Mange kunde være overbeviste om dens Rigtighed, hvis de ikke for Beqvemmeligheds og Fordeels Skyld bedroge sig selv.

Ogsaa Svaghed kan være Affectationens Kilde. Menneskene kunne af Velvillie saaledes lade sig henrive, at de paatvinge sig selv en Overbeviisning for at være enige med den, de elske. Hvorledes kunde det ellers gaae til, at Folk under videnskabelige Feider, hvor Overbeviisningen blot skulde bestemmes af Grunde, dog i Almindelighed ere paa samme Partie, som deres Venner, Frænder og Paarørende? \opage{175}Men denne Lyst til at være enig med Andre og derfor at blive det, kommer ikke altid af Velvillie: den kan ogsaa have sin Grund i Forfængelighed. Dette er den anden Svaghed, som tidt foranlediger Affectation. For at være behagelig for Andre bliver En tidt enig med dem. Dette Slags Affectation findes især hos sangvinske Mennesker og kan, naar de have udmærkede Aandsgaver, frembringe mærkelige Phænomener. Der gives Mennesker med særdeles hurtige Hoveder, som i den Grad ere et Product af deres Omgivelse, at de gandske antage dens Farve. Høi Grad af Smidighed i Aandens Evner og Receptivitet for de forskjelligste Synsmaader gjøre, at de snart slutte sig til eet, snart til et andet literarisk eller politisk Partie, og ere dog med Liv og Sjæl paa ethvert af dem, saa længe det varer. Saadanne usammenhængende Mennesker leve nemlig altid i en umiddelbar Erkjendelse, uden at kunne hæve sig til det Standpunct, hvorfra forskjellige eensidige Standpuncter kunne overskues og sees i deres relative Gyldighed. Naar dette Slags Vendekaaber besidde høi Grad af Aandsrigdom, som de ofte gjøre, kunne de paa en glimrende Maade udføre den Person, der er dem givet udenfra, uden at have nogen dyb Rod i deres egen Personlighed.

Det er ikke blot en forfængelig Lyst til at være med og finde godt Kammeratskab, der er Skyld i saadan Nærlighed, det kan ogsaa være en Stræben efter den Nydelse, som forhøiet Selvfølelse skjenker. (\emph{Moralske Drømmerier}).\footnote{\textbf{Alibi:} Exempel paa den af Lyst til at behage sig selv udspringende Affectation er Phantasier over Planer til Virksomhed, som man halv om halv veed, man ei vil udføre.}

Enhver Personlighed er noget uendeligt Herligt, naar den udføres efter sin Idee; naar man derimod lader den overvældes af fremmede Hensyn, forfeiler man sin Bestemmelse. Maaskee \opage{176}mere farlig for Moralitet end bevidst Løgn er den med Selvbedrag blandede Usandhed, hvorved Mennesket amalgamerer sin falske Person med den virkelige. Den fuldbevidste Løgn kan man aflægge efter et alvorligt Forsæt; men man kan voxe fast i den falske Person uden at kunne rive sig løs derfra. Man gjør sig selv til en moralsk Skifting.

\ombreak

Naar man overlader sig trøstigt til en Idee og følger dens Veiledning uden ængstelig at forsone hver Sætning med gjængse Synsmaader, da ansees man gjerne for phantastisk. Det vil jeg vist blive anseet for, naar jeg i mit Værk om Affectation paastaaer, \emph{at ingen Livsyttring har Sandhed, uden deri ligger skabende Selvvirksomhed} (S.\ 91).

I Indledningen (om Begrebernes Tyrannie) maa vel udvikles, at man ei tænker i Ord blot.

Der gives to Slags Sandhed i Menneskelivet. Den ene er naturlig, medfødt: saaledes kan Tilbøielighed til Affectation være medfødt. Den anden er med Frihed valgt og tilvundet; den er af høiere og rigere Indhold.

Noget Selvbedrag er nødvendigt; thi det udgjør Livet.

\ombreak

Affectation kommer oftest af, at man ei har Kraft til at lægge sig ud med Verden for at vise sin Personlighed sand. Derfor er det godt, at Somme komme til at staae i trodsig Opposition til Lauget. I Naturstanden levede hver adskilt og udviklede sig trodsig en Personlighed, da han ei forstyrredes af Mange. Nu danner man sig ved Abstraction en almeengyldig \opage{177}Person, et Selskabsideal uden Kanter, et Ideal uden Individualitet. (Naturen gjør da sin Ret gjeldende i Bøger som Baggesens Gjenganger.)

Ogsaa Opposition kan drives med Affectation; man kan forstærke Disharmonien, fordi man har offret sit Liv til Krig.

Enhver har af Naturen sit dybe Præg, men lader det udslettes for at tækkes Andre. Disse Andre have i at forkaste de Fremstikkende et meget rigtigt Princip; kuns skulde de gjøre Forskjel imellem de Selvraadige og dem, der ere fremstikkende af Pedanterie og Hjelpeløshed.

De aldeles opløste Mennesker, hvis Liv kun er Mimik, have undertiden saa riig en Natur, at de kunne udføre en paataget Charakteer med en saadan Kraft, at de Individer, som virkelig have denne Charakteer, ei kunne gjennemføre den kraftigere.

De Usammenhængende ere for slette til at være onde.

\ombreak

Affectation kan komme af, at man bilder sig ind, at man virker for en Idee, og skaffer sig selv derved en forhøiet Selvfølelse.

Somme tilskrive sig selv falske Egenskaber blot for at sige Noget om sig selv. --- Hysteriske Qvinder ---.

Affectation af Naturfortrin: Naar en Ludimagister vil være en Ungdommens Gud-Apollo. Eller naar man bander og bruger aposterioriske Vittigheder, for at synes en Kraftaand, da det dog egentlig blot er Magelighed, der bringer En til Sligt. Eller \opage{178}naar en gammel Jomfru vil være ung. Eller naar en Doctor Faust vil være romantisk Elsker, eller en gammel Gjæk være fidel med unge Studenter; Rahbek være kjæk Kriger. --- Mellem Leie og Eie svæver Livet. Ungdommen raser, sagde Kjærlingen, hun sprang over et Halmstraa. Cicero vil være en Frihedshelt. Ved hans Frihedsiver udrettedes Intet; den var blot til Stads.

Man skal netop blotte store Mænds Feil, for ret at lære den menneskelige Natur at kjende. Næsten alle dannede Folk lyve i Fortællingen om deres Kjærlighed.

Det er ofte en Feiltagelse, at Mangen troer sig i Stand til at skrive store Værker; han holder sig selv en falsk Bebudelsesfest.

At kokettere med Aabenhjertighed, Beskedenhed o.\ s.\ v.; med bevidstløs Overgivenhed, Naivetet, barnligt Sind. Den affecteret Trohjertige og Aabenhjertede kniber sit Øie lidt sammen; Selvtilfredsheden, der lyser af hans Øine, forraader ham, og det halvbetvungne Smiil, der spiller om hans Mund.\footnote{Sml.\ S.\ 54.}

„Asnet kunde ikke see, det netop var barnligt af mig, at jeg gjerne vilde æde Pandekager.``

En af Aarsagerne til, at menneskelige Bevæggrunde ligge saa gandske i Mørke, er den Omstændighed, at Folk, for at synes ædle, sige, ved at høre Observationer om Egoismen: Nei, Gud bevares, det er ikke sandt. --- Ogsaa ved at misbillige den mindste Plet paa Tragoediehelte lægge de deres ædle Sjæl for Dagen. Derimod Menneskefjender overdrive Alt. --- (Spleenfulde Englændere eller franske fortvivlede Gudsfornegtere). \opage{179}De, der pynte sig med Moralitet (hvilket er ligesaa galt som at pynte sig med Immoralitet), bruge: „Gud forlade ham``, i samme Betydning og med samme forbittrede Blik, som naar de vilde sige: „Gud forbande ham.``

Tjenestepiger sige for at tirre hinanden: Gud forlade dig din Synd, Grete! men du er en ond Satan --- i en sagtmodig Tone, saa at de komme til selv at troe sig ædle.

(At omgive sig med en Helgenglorie).

At Ingen let for en Søster eller Broder forfalder til moralske Declamationer, viser, at den rigoristiske Declamator bruger sin Iver som en Stadsdragt.

Hvor stor er ikke Forskjellen paa fjortenaarige Pigers og Drenges Opførsel i Fremmedes Nærværelse? Disse skjendes uartigen med deres Moder i Fremmedes Paahør, hine opføre sig gracieuse --- af en ubevidst Lyst til at behage.

Pigen kjæler for Katten for at synes ømtfølende.

Det er tidt for at give sig Skin af Alvidenhed, at Somme sige: Alt er gammelt.

En Affectation hos lærde Mænd og Genier er det, at synes ei at være af denne Verden. De høre med Fryd deres Koner tale om deres Distraction, Sorgløshed og Ukyndighed i det daglige Livs Anliggender, da de ei ere af denne Verden.

Betænk, hvormegen Løgn der kan findes i en dagligdags Scene af Livet. Den Reisende spørger Bernhard: „Hvorledes \opage{180}laver man Pandekager?`` Lovise: „Det veed min lærde Broder ikke.`` Bernhard: „Aa, jo vist veed jeg`` (opgiver paa Skrømt nogle urigtige Elementer og støder sin Naboerske Augusta i Siden). Hun siger: „Ja med Skam at tale om`` (mener hun det?) „veed jeg det heller ikke`` (siger hun deri sandt?).

Undertiden spiller en heel Stand sin Rolle som S.\ og fornemme Matroser.

Matroser af Stand mumle efter de egentlige saakaldte Ulke med halvaabnet Mund eller bruge et Mæle, som om de havde Skraa i Munden, havende ingen.

\ombreak

Det er en snurrig Slags Svaghed hos visse Folk, at Høfligheden tyranniserer dem. --- \textbf{A.} „De maa endelig drikke Kaffe hos Mini, den er bedre end Semadenis.`` \textbf{B.} „Nei; jeg gaaer heller til Semadeni.`` --- \textbf{A.} „Jeg forsikkrer Dem, De maa absolut drikke Deres Kaffe hos Mini.`` --- \textbf{B.} „Nei`` --- \textbf{A.} (skydende ham ind) „Jo, det maa De absolut.`` --- \textbf{B.} lader som han vil derind, og lurer i Porten, til \textbf{A.} er noget borte; derpaa gaaer han sin Vei. Den Enes paatrængende Høflighed og den Andens nathueagtige Delicatesse er lige mærkelig. Saa rørte og betvungne blive Folk af Høflighed og Omsorg for dem, som de paa begge Sider vide er et Skin.

Det er et meget uskyldigt Slags Affectation, at sætte Musklerne i den Stilling, som de indtage, naar man leer eller smiler, imens man hører en Anekdot, som ikke morer, men dog gjør Fordring derpaa. \opage{181}Den, der glæder sig over Complimenter ved Disputatser, veed dog meget vel, at de, der sige dem, ei mene dem. Han bedrager frivillig sig selv for at nyde en liden Triumph. --- De Store vedblive ogsaa at troe paa Smiger.

Roes og Beundring blive med Tiden til en physisk Nødvendighed for gamle Mænd. De ville leve i en Drøm om at have været Noget, uagtet de Intet ere; ville bedrage sig selv.

Forsoning er tidt Selvbedrag. „Nu har jeg sluttet Forlig med ham, den Æsel!``

Er det ei Affectation, naar en stolt Choleriker vil skjule, at han forundrer sig ved Noget eller finder det latterligt \&c., men yttrer, at han ei finder det saa?

Affectation reiser sig undertiden af Hovmod. Den, der vil skjule den Virkning, Andres Tale gjør paa ham, er paa eengang usand og stedt i Selvbedrag. Det Sidste, fordi han ikke i samme Øieblik, som Usandheden yttres, er sig den bevidst, skjøndt han ved en følgende Reflexion tidt maa kunne erkjende sin Falskhed.

\emph{At skabe sig til.}

Skræk for Affectation kan --- saa besynderligt det end ved første Øiekast synes --- drives saavidt, at en abnorm Sindstilstand derved fremkommer. Denne Skræk frembringer en misforstaaet Holden paa den abstracte Identitet med sig selv. Saaledes har jeg kjendt en ung Mand, som ikke kunde overtale sig til at bortsende noget Brev, der havde ligget en Dag over i hans Pult, fordi det forekom ham at være usandt, da det nemlig nu ikke meer udtrykte hans nærværende Sindsstemning.

\ombreak

\opage{182}En Affectation er livløs Eftersnakken af de philosophiske Folkesagn om Kunstens Herlighed. Ei blot Kunst i ængere Betydning, men Livets Kunst fortjener Priis. --- Er Besindigheden ei Enthusiasme, da er der ingen Enthusiasme i Naturens rolige fornuftige Skaben, og det er jo dog just den, der skal efterlignes. --- Ørhed --- Blodet til Hovedet --- drukken Forrykthed, og den saakaldte Begeistring, som egentlig er en forrykt Glæde over Blodets Syden.\footnote{Den momentane Exaltation er tidt af en lav sandselig Natur. Forf.\ Anm.\ Jvfr.\ S.\ 89, tredie Aphorisme.} --- Det er ofte Forfængelighedens Henrykkelse over enkelte heldige Udtryk. Enthusiasmen er den rolige Arbeiden af et af Naturen i os nedlagt Mønster. Det kan bedre efterlignes ved rolig Flid, end ved Susen for Ørene. Der skal kun være Nydelse i Følelsen af det Heles Lykken, ei i det enkelte Ord og Udtryk. Forudsætter en rolig Udmeislen af et græsk Billed ei Enthusiasme? Det er en altid tilbagevendende Affectation, at skjule sin Flid.

Alt i Homers Tid blev Talen behandlet som et Kunststykke. --- Homer er ei skrevet i Ruus.

Løgn er den Poesie, som ei kommer af Livet. Jo nærmere den kommer Livet, jo mere sand; jo fjernere, jo mere Løgn.

Den moralske Phantast mangler Phantasie; thi hans Phantasmer ere abstracte.

Den nærværende Generation har en saa høi Grad af Incitabilitet, at den bliver sat i Sindsbevægelse ved at bevise Sindsbevægelsens Mislighed med Hensyn til kunstnerisk eller videnskabelig Production. Ved den levende Forestilling om den Høihed, der ligger i Erkjendelsens og Villiens Overvægt, og den usunde \opage{183}Retning, der lægges for Dagen i det gjængse Føleri, finder den Næring for sit eget Føleri. De meest i Følelse Opløste, som nødvendigen maatte komme til Bevidsthed om dette deres Føleri, stride med den største Pathos imod Føleriet.

\ombreak

Villien skal ei bestemme, hvad der er sandt.

Uægte Enthusiasme kan yttre sig deri, at man ved heftig Declamation vil gjøre sin Villie gjeldende, uden at tilveiebringe Overbeviisning ved Grunde.

Mange ville ikke af med deres Vildfarelser, der med lange forviklede Rødder have som haardnakket Ukrud snøret sig fast om særdeles mange Naboplanter. De troe at stride for Sandhed og stride for deres egne Personer.

Det er altid Affectation, naar en Dom ikke udtaler en Tanke, der er bestemt ved et Object, men tildeels ved en Tilbøielighed. Den, der har i Sinde at skrive, betragter sig som hørende til Skribenternes Partie og er tidt overbærende, men bliver overdreven streng, naar han har resigneret.

Folk lukke Øinene for Grækernes Skyggeside og Tyranners Lysside.\footnote{See S.\ 37.} Affectation er det, naar En, der i lang Tid har viist Tilbøielighed til at fremhæve Skyggesiden ved Alt, pludselig vil vise sig villig til at erkjende den bedre Side ved Tingene og pludselig giver sig til at rose Et og Andet. --- Ligeledes naar Folk af Frygt for at vise Følsomhed anstille sig haardhjertede, --- og \opage{184}naar man tvinger sig til at erkjende fornemme Folks Fortjenester, for ei i sine egne og Andres Øine at gaae for en Smigrer.

Et Slags Affectation er den uphilosophiske Mængdes foragtelige Yttringer om en Philosophie, der er over deres Horizont. De sætte deres egen Erkjendelses Grændser, men med svag Overbeviisning, som Fornuftens Grændser. (Den, der ikke kan forstaae Philosophien, vil helst, at Intet skal være i den.)

\textbf{Par force} at ville ynde noget Snavs af Frygt for sin egen Criticismus. --- Det er en Affectation, at fornegte sin kritiske Natur.

Naar Drengen siger: jeg gjør mere af Goethe end af Schiller, eller overhovedet den Umodne gjør mere af den objective end den subjective Poesie, da er det, som naar Barnet negter sig at elske Sødt, og tvinger sig til før Tiden at elske det Bittre.

Affectation er Tegnérs Tale om Forms og Stofs Sammensmeltning hos Grækerne, hvilken han selv modsiger ved at levere overveiende Form.

\ombreak

Siger Den Sandhed, der taler Ord, han ei selv forstaaer? --- Nei.\footnote{Jvfr.\ S.\ 30.} Siger da Den Sandhed, der betjener sig af Ord, hvis Begreb er ham bibragt af en Anden, uden at de staae i nogen organisk Sammenhæng med hans øvrige Erkjendelse? Taler den Sandhed, der disputerer med Ord, han ikke forstaaer, der beherskes af sine egne Definitioner? --- Men, det er jo blot \opage{185}Ordspil; den, der troer, at han taler Sandhed, er \textit{bona fide}. --- Nei, han bedrager sig selv. Han skammer sig ved at staae nøgen, idet han bortkaster alle traditionelle Begreber, som udgjøre hans hele videnskabelige Forraad. Han vil ei komme frem med sin egen Person, troer ei paa sin Persons uendelige Dybde. Naar hver enkelt Mand, uden at befrygte Dadel for Eenfoldighed, dømte om Tingene saaledes, som de fremstillede sig for ham, da vilde herlige Phænomener fremkomme. Men de Fleste frygte for at lægge Aands-Armod for Dagen. De bestyrke forvirrede Hoveder i deres Forvirring, ved at pine sig ind i disses Synsmaader, af Frygt for at ansees for Folk af ringe Fatteevner. --- Derfor pine danske Folk sig til at tale germanisk; derfor agte Folk philosophiske Phraser, som de ei forstaae.

Den taler ikke Sandhed, som siger, at Mennesket ei giver sig selv Love; thi han veed ikke, hvad et Jeg er.

En Affectation er Moralitet efter døde Begreber. Den, der f.\ Ex.\ siger: „Fy! A er en Selvmorder``, har kun et Ord at holde sig til. Den, der siger: „A burde ikke forlade sin gamle Fader``, han taler af Anskuelse, umiddelbar Erkjendelse.

Et eget Slags Affectation er det af en Eklektiker, naar han bruger systematiske Philosophers Terminer; hvilket er unødvendigt for ham, hvis uorganiske Philosophie kan forstaaes paa ethvert Punct.

Folk holde sig af Fornemhed i de generelle Begreber.

\ombreak

Et eget Slags Affectation i Foredraget er den grammaticalske. Denne foraarsager, at lexicalske Studier og Omhu for \opage{186}en skoleret Syntaxis træder synlig frem i Stilen. Naar man læser slige Skrifter, er det, som man læste et dødt Sprog. Deslige Purister anlægge en Fabrik for indenlandske Ord, som skulle afløse andre gangbare, hvis Indfødsret er mistænkt. Disse hjemmegjorte Ord gjemme de i deres Sparebøsse, og naar de have sanket et passende Antal, tage de deraf Anledning til at skrive over en eller anden triviel Gjenstand, som tillader dem at skjenke Ordbygningen udeelt Opmærksomhed. Naar man ret betragter deres Producter, er Sproget det egentlige Indhold, og det tilsyneladende Indhold en Biting.

Affectation: at skrive i en afgjørende Tone, voldsomt at slutte sin Erkjendelses Regnskabsbog. Det tvinger man Børn til ved at lade dem skrive danske Stile. Man gjør dem til Stylistikere, som er noget af det Værste, et Menneske kan være.

Det er underligt, at Folk ikke tabe Lyst til at være Stiilskrivere ved at see, med hvilken Lede de læse Andres Stile. Egentlig skulde Stile blot skrives med negativt Velbehag, med den Overbeviisning, at man brugte passende Udtryk.

Glæde ved at bruge visse Ord: Dannemand, hartad.

Et Beviis paa, at Glæde over Stiil er Løgn, er at en Dreng rødmer, naar man gjør ham opmærksom paa, at han bruger dvæle for tøve, gjenmæle for svare, eller declamerer om sølvhvide Oldinge.\footnote{Om Stylisteri jvfr.\ S.\ 131, 35, 23.}

\ombreak

\opage{187}Der var engang en Mand, som sjelden gik ud at spadsere, uagtet hans Legemsbeskaffenhed gjorde det ønskeligt, at han havde gjort det. Dette Savn af Bevægelse blev derimod ved et Naturinstinct erstattet paa en anden Maade; han talte sig nemlig i Affect over en eller anden Materie, og hans Blod sattes derved i en saadan Circulation, at han ei havde anden Motion behov. Da enhver Sag kan sees fra to Sider, var det tilfældigt, fra hvilken Side den øieblikkeligen blev seet. Det øieblikkelige klare Blik for en isoleret Sætnings Rigtighed, taget i den subjective Mening, som han just tilfældigviis tog den, tjente ham til Vehikel for en Tummel i Forestillingskredsen, der snart igjen lagde sig, uden at efterlade mærkelige Følger deri. Forfatteren af denne Bog har engang i China seet et Knippe Frøer, der holdtes tilfals paa et Torv, og som i Almindelighed laae uden mærkelig Bevægelse, fordi de sammenknyttede Dyr anstrengte deres Kræfter i saa modsatte Retninger, at Resultatet af deres Anstrengelser omtrent blev Nul. Undertiden lykkedes det derimod en enkelt Frø uden Modstand ved et dristigt Hop at løfte hele Knippet for et Øieblik i Veiret. Saaledes gik det med den omtalte Mands Forestillinger; undertiden blev en enkelt Forestilling saaledes herskende, at hans Foredrag antog et forbigaaende Skin af systematisk Sammenhæng.

Den videnskabelige Arbeidsmand samler mechanisk, hvad der pleier at staae sammen i en Disciplin, uden Erkjendelse af dets Sammenhøren.

En Affectation er det, at den, der vil skrive en Afhandling, tænker sig en Titel og beskrænkes deraf.

Det er et Slags Affectation, at gjøre en Afhandling fuldstændig, fordi enkelte Capitler ere levende for os. \opage{188}At de Gamle ei havde den systematiske Affectation, sees af deres hyppige Svinkeærinder. De tyranniseredes ikke af deres Titler. Alt, hvad der er Liv, strider mod stive Former.

Affectation er Symmetrie i philosophiske Systemer.

\ombreak

Det fremmer Menneskekjærlighed, at tydeliggjøre sig Begreber om slette Egenskaber, for at de kunne skarpt adskilles fra de gode. (Somme finde Egoismus i Alt.)

En begyndende Menneskekjenders Indsigt i Slethed er Selvbedrag. Vogt dig at komme i hans Nærhed. Han begynder, som Rollingen, der ønsker at lære at kjende en Ting, med at rive itu. Men al Formodning om Slethed eller Herlighed svinder i Nærheden.

Utaalelige Menneskekjendere, som ville tvinge Folk til at handle efter deres Forestillinger, og Usselrygge, som ikke tør sætte sig imod hines Begreber!\footnote{Jvfr.\ S.\ 21, 26.}

\ombreak

\end{document}
